% !TeX root = ../papers/nc-p3.tex
\section{非対称成分の滑らかな種数三切断による復元}
\label{sec:finite-genus-three-curve-cut}
\begin{theorem}\label{thm:finite-genus-three-curve-cut}
定理\ref{thm:quartic-center-order}の十五次元成分の開集合にある任意の窓は、
滑らかな平面四次曲線 $C$ と四つの半安定な階数4の束 $E_0,\ldots,E_3$ を用いて復元できる。
$E_{i+4}=E_i\otimes\omega_C$ と置けば $\deg E_{i+1}-\deg E_i=4$ であり、
一般の二本の中心四次式 $g_1,g_2$ を割った全経路代数は
\[
 D_{i,i+n}=H^0(C,E_i^*\otimes E_{i+n})\quad(n\ge0)
\]
の合成代数である。特に二つの隣接窓は $n=2$ の積の核として定まる。
三次元射影中心が二十の既知節点以外で滑らかであるという仮定を置く必要はない。
\end{theorem}
\begin{proof}
五変数の射影四次中心は整かつ正規なので、その特異集合は次元1以下である。
一般平面はそれを避け、Bertiniにより滑らかな整平面四次曲線を切る。
Azumaya軌跡の命題\ref{prop:smooth-azumaya}により全曲線上にAzumaya代数を持つ。
一般の二式は四つのNoether正規化の一次部分の中の最初の二元にできる。
実際、四つの一次式で割ると有限になる条件は開で非空であり、一般の二式に残り二式を足せばよい。
中心と折り畳み環は次元4の極大Cohen--Macaulay加群であるから一般の四式上で自由である。
あるいは一組の既知の多項式自由表示から深さ4、従って任意の斉次パラメータ系の正則性を得る。
二式を割った隅は残りの二変数上の自由加群で、四移動ごとにまとめても深さ2を持つ。
よってこれらは射影曲線上の対応する層の全切断加群に等しい
（無関係イデアルの零次・一次局所コホモロジーが消える）。

一般平面上でこの層表示を具体化する。経路長を $4m+j-i$ にまとめた $ij$ 隅を次数 $m$ として層化する。
中心次数1が $\mathcal O_C(1)=\omega_C$ の捩りに対応する。
全隅を合わせると一般中心点で四冪等元を持つ次数16のAzumaya代数である。
曲線の各点で可逆な中心一次式により剰余次数を自明化して見ると、
これは元の次数環をその値1で割った局所表示である。四射影の行列階数は4ずつで一定となる。
幾何学的曲線上でBrauer群は零なので、全体を一つの束の自己準同型として分裂できる。
積の向きを合わせて必要なら反対代数を取り、その四成分を $E_i$ とすれば上式を得る。
二式の正則性より各始点のHilbert級数は $(1-z^4)^2/(1-z)^4$。
大きな $m$ でRiemann--Rochの $4\,(\deg E_j-\deg E_i)+64m-32$
をこれに比較すると $\deg E_j-\deg E_i=4(j-i)$ となる。

半安定性は固定窓の一般平面で示せる。平面の関数体上の曲線を取ると、
その関数体は射影中心の関数体の純超越拡大である。実際、点とそれを含む平面の入射の射影中心上の
一般ファイバーは点の商空間における二平面のGrassmann多様体で有理である。
従って系\ref{cor:quartic-division-corners}の次数4可除性は保たれる。
四つの束をこの体上でねじれ束として読むと、ねじれた一般階数は4の倍数に限られる。
幾何学的不安定部分は降下して階数1から3の部分を与えてしまうため存在しない。
一般の閉平面での半安定性は開性で従い、Brauer分裂の共通線束の自由度にも依存しない。
\end{proof}

\begin{proposition}\label{prop:finite-genus-three-curve-recognition}
逆に滑らかな種数3の $C$ と上記の次数・半安定性の四束を与える。
$n=0,1,2,3$ の切断次元が各始点で $1,4,10,20$ であるとし、
四つの一次空間の積から四つの二次核を取り、それらで定義した折り畳み環を $B$ とする。
次の範囲の条件だけで、$B$ が標準左右AS正則であり、二次元の中心四次鉛筆 $G$ を割ると
ちょうどこの切断代数になるような表示を判定できる。
\begin{enumerate}
\item 標準有限Frobenius条件と、曲線への像が零である独立な四次中心 $g_1,g_2$ を検査する。
\item $B/(g_1,g_2)$ から曲線の切断代数への評価が経路長13まで全端点で同型であることを課す。
\item 元の経路次元を36次、標準複体の完全性を16次まで調べる。
\end{enumerate}
元の窓が直前の定理の開集合に属することを逆判定の仮定に含めない。
$G$ を指定した折り畳み環（次数・頂点を固定）の同型は、曲線の同型と四束の共通線束捩りまでの同型と必要十分である。
鉛筆を忘れた添字保持同型は、復元した二窓の三基底変更で判定できる。
\end{proposition}
\begin{proof}
半安定性から $n<0$ の切断は零で、$n>4$ ではSerre双対により第一コホモロジーが零となる。
$n=4$ では第一コホモロジーは自己準同型の双対で一次元である。
Riemann--Rochと仮定から全切断代数のHilbert級数は先ほどの二式の商に一致する。
標準束に共通零点のない二切断 $s,t$ を選ぶ。各隅でこれらは切断加群の正則列である。
実際、束の無捩れ性により $s$ は非零因子で、$tf=sf'$ なら $f/s$ は $s$ の零点でも正則である。
負の長さの切断は零なので切り詰めによる例外も生じない。よって二式を割ると
各始点で $(1+z+z^2+z^3)^4$ が残り13次以後は零である。
この二切断の恒等倍を全頂点で取り、第二検査により $B/(g_1,g_2)$ へ持ち上げる。
五次の評価が単射なので全矢とそこで可換である。これで前定理の順次中心四元が揃う。
十三次までの同型より最終商の必要次元を得、第三検査と定理\ref{thm:successively-central-cutoff}が
全次数正則性と四元の順次正則性を与える。
切断代数は $s,t$ と十二次以下の切断で生成されるので評価は全次数で全射であり、
二中心四次の正則性からHilbert級数も一致し同型となる。
逆にこの表示があって正則なら四条件の次元・有限検査は自動的に満たされる。

最後に $G$ を割った零番隅の四Veronese環の中心は $\bigoplus_m H^0(m\omega_C)$ である。
実際高い捩りの全切断が各点で自己準同型束を生成するので、中心元は一般点でスカラー、
そのスカラー切断は全点で正則である。この中心のProjと偏極から曲線と標準束を取り戻す。
隅を層化すると $\mathcal End(E_0)$ が得られ、その分裂束は線束捩りまで一意である。
次に他の三隅の $\mathcal Hom(E_0,E_i)$ を加群として読めば同じ共通捩りで $E_i$ が定まる。
逆向きにはこの曲線データから積と二次関係を復元できる。
鉛筆も評価の四次核のうちの頂点整合中心として一意である。実際 $B/(G)$ で消える四次元は
各頂点で二本の独立な $g_1,g_2$ の成分であり、隣り合う頂点の係数差が作る一次関係は
二元が正則列で全次数商のHilbert級数が上の通りであることから零となる。
具体的には二元の最初の係数シジジーは全次数8で、次数5にはない。
従って全四頂点で同じ二係数を持たねばならず、消える中心は正確に $G$ である。
\end{proof}
